\documentclass[11pt]{article}
\usepackage[T1]{fontenc}
\usepackage[utf8]{inputenc}
\usepackage{amsmath,amsthm,mathtools}
\usepackage{newpxtext,newpxmath}
\usepackage[a4paper,margin=28mm,headheight=15pt]{geometry}
\usepackage{microtype}
\usepackage{enumitem,booktabs,longtable,array}
\usepackage[dvipsnames]{xcolor}
\usepackage{fancyhdr}
\usepackage[most]{tcolorbox}
\usepackage{xurl}
\usepackage[colorlinks=true,linkcolor=MidnightBlue,citecolor=MidnightBlue,
 urlcolor=MidnightBlue,pdftitle={Polish Presburger Arithmetic inside the Omnific Integers},
 pdfauthor={Research manuscript prepared with AI assistance for Vladimir Reshetnikov}]{hyperref}
\usepackage[nameinlink,noabbrev]{cleveref}
\usepackage{listings}
\definecolor{ink}{HTML}{15394B}
\definecolor{paperblue}{HTML}{F0F5F7}
\tcbset{colback=paperblue,colframe=ink,boxrule=.45pt,arc=1mm,
 left=3mm,right=3mm,top=2mm,bottom=2mm}
\lstset{basicstyle=\ttfamily\small,breaklines=true,columns=fullflexible,
 frame=single,rulecolor=\color{black!20},aboveskip=1em,belowskip=1em}
\setlist[enumerate]{itemsep=.35em,topsep=.5em}
\setlist[itemize]{itemsep=.3em,topsep=.5em}
\setlength{\parskip}{.35em}
\setlength{\parindent}{1.2em}
\setlength{\emergencystretch}{2em}
\pagestyle{fancy}
\fancyhf{}
\fancyhead[L]{\small\scshape Polish Presburger Arithmetic}
\fancyhead[R]{\small Research manuscript}
\fancyfoot[C]{\thepage}
\renewcommand{\headrulewidth}{.3pt}
\numberwithin{equation}{section}
\newtheorem{theorem}{Theorem}[section]
\newtheorem{lemma}[theorem]{Lemma}
\newtheorem{proposition}[theorem]{Proposition}
\newtheorem{corollary}[theorem]{Corollary}
\theoremstyle{definition}
\newtheorem{definition}[theorem]{Definition}
\newtheorem{example}[theorem]{Example}
\newtheorem{question}{Research question}
\theoremstyle{remark}
\newtheorem{remark}[theorem]{Remark}
\newcommand{\N}{\mathbb N}
\newcommand{\Z}{\mathbb Z}
\newcommand{\Q}{\mathbb Q}
\newcommand{\R}{\mathbb R}
\newcommand{\No}{\mathbf{No}}
\newcommand{\Oz}{\mathbf{Oz}}
\newcommand{\Baire}{\N^{\N}}
\newcommand{\lex}{\mathbin{\overrightarrow\times}}
\newcommand{\MR}{M_{\R}}
\newcommand{\MB}{M_{\mathrm B}}
\newcommand{\GD}{G_D}
\newcommand{\MD}{M_D}
\newcommand{\res}{\operatorname{res}}
\newcommand{\Th}{\operatorname{Th}}
\newcommand{\Aut}{\operatorname{Aut}}
\newcommand{\End}{\operatorname{End}}
\newcommand{\id}{\operatorname{id}}
\newcommand{\ct}{\operatorname{ct}}
\newcommand{\one}{\mathbf 1}
\newcommand{\Dtwo}{\boldsymbol{\Delta}^{0}_{2}}
\newcommand{\proofstatus}{\textit{Proof status.}}
\title{\vspace{-1.5em}\color{ink}\bfseries
Polish Presburger Arithmetic\\[.3em]
inside the Omnific Integers\\[.6em]
\large An affirmative construction for Glazer's Question 2,\\
effective Baire-space models, and structural obstructions}
\author{\normalsize Research manuscript prepared for Vladimir Reshetnikov}
\date{October 3, 2026}

\begin{document}
\maketitle
\thispagestyle{empty}
\begin{abstract}
We construct an uncountable Polish model of the full first-order theory
$\Th(\N;0,1,+,<)$ whose addition is continuous. The simplest carrier is
\[
 \MR=(\R_{>0}\times\Z)\ \cup\ (\{0\}\times\N),
\]
with coordinatewise addition, lexicographic arithmetic order, and the
subspace topology from $\R_{\rm usual}\times\Z_{\rm discrete}$.
This gives an affirmative answer to Question 2 as stated on page 8 of
Elliot Glazer's \emph{A Topological Tennenbaum Theorem}
(arXiv:2311.13699v1). A second construction, using the lexicographically
ordered group $\Q^{\N}$ with discrete rational coordinates, is explicitly
homeomorphic to Baire space and has Type-2 computable addition. We prove
quantifier elimination over the relevant ordered groups, including the
full induction scheme on their positive cones. All definable relations
are $\Dtwo$; in the Baire presentation they admit uniform finite-mind-change
truth approximations. Rank two is optimal for the order relation in any
uncountable Polish continuous-addition model.

Both constructions are set-sized additive fragments of Conway's omnific
integers. They admit no unital semiring multiplication whatsoever with
the displayed addition and unit. We also classify the real model's
unital endomorphisms and elementary submodels, exhibit proper clopen
elementary self-copies of the Baire model, and isolate the stronger
coset-invariance hypothesis needed to extend ProveIt's integer-specific
Cooper lemmas to nonstandard groups.
\end{abstract}

\begin{tcolorbox}[title=Status and scope]
The mathematical assertions designated as theorems are accompanied by
proofs. The construction answers the cited question \emph{as written};
independent bibliographic priority has not been established. Classical
Presburger elimination and elementary topological-group facts are not
claimed as new. This AI-assisted manuscript is not independently refereed
or Lean-verified. The accompanying executable tests are finite checks, not
a replacement for the proofs.
\end{tcolorbox}

\clearpage
{\setlength{\parskip}{0pt}\small
\tableofcontents
}
\clearpage

\section{The research overlap and the precise target}
\label{sec:overlap}

The relevant intersection is \emph{topological and effective model theory
of weak arithmetic}. Glazer's work supplies a sharply formulated question;
ProveIt supplies two concrete mathematical interfaces: an executable
Cooper quantifier-elimination development and a normal-form construction
of the omnific integers. This makes it possible to connect an external
Polish topology, internal induction, and exact surreal coordinates without
appealing to the repository's unreviewed research claims.

In the final section of \cite{Glazer}, Glazer asks:
\begin{quote}
Does some uncountable Polish space support a model of Presburger
arithmetic with continuous addition?
\end{quote}
The source is Question 2 on page 8 of arXiv:2311.13699v1. We use precisely
that formulation. The topology need not be the arithmetic order topology,
and no multiplication is required. These two distinctions are essential.

Glazer's obstruction for stronger arithmetic includes the impossibility
of an uncountable Polish model of $Q+I\mathrm{Open}$ with continuous
addition and Borel multiplication \cite[Corollary 1]{Glazer}. The earlier
study of Enayat, Hamkins, and Wcis\l o \cite{EHW} motivates the broader
problem of finding topological presentations of arithmetic. Our
construction does not contradict those results: it supports the additive
language, and its additive structure prevents a semiring expansion.

\subsection{Repository interfaces actually inspected}

The repository snapshot used here is
\begin{center}
\texttt{a11efab09a2d08f692a09829cb6d3440ddd1795a}.
\end{center}
The Presburger README \cite{RepoPresburger} describes a complete Lean
integer quantifier eliminator and an independent Rocq/Coq proof of the
normalized one-variable step. The inspected file
\texttt{PresburgerArithmetic/Cooper.lean} \cite{RepoCooper} states
\texttt{periodic\_interval} and \texttt{periodic\_has\_residue} for
predicates on the ordinary type \texttt{Int}. Their hypotheses are entirely
appropriate there. They must not be transferred unchanged to a
nonstandard ordered group; \cref{sec:periodicity} gives an explicit
counterexample and the correct replacement.

The surreal README \cite{RepoSurreal} records an omnific ring whose normal
forms have nonnegative exponents and integer constant term. Only this
normal-form interface and the standard arithmetic of its additive group
are needed below. We do not assume that the present topology or any of
this article's theorems has already been formalized in that repository.

\subsection{What is established, and what is not}

The principal result is an explicit affirmative construction, rather than
a claim that a difficult negative theorem has been overturned. The
lexicographic algebra is classical. The contribution developed here is
the precise topology and its application to the stated question, together
with an effective Baire presentation and a collection of sharp structural
consequences.

The literature search verified the question in the source PDF and checked
the arXiv record; it did not establish the absence of an earlier solution.
Consequently, ``an affirmative answer to the stated question'' is the
appropriate conclusion, not an independently certified first solution.
No answer to Glazer's separate reverse-mathematical Question 1 is claimed.

All constructions and proofs take place in ordinary ZFC with set-sized
carriers. The notation $\Oz$ is used only to identify where explicit
set-sized images lie. No topology on an unrestricted proper class, and no
set of all surreal numbers, is used.

\section{The construction in one theorem}
\label{sec:main}

We include $0$ in $\N$. For an ordered abelian group $D$, let
\[
 \GD=D\lex\Z
\]
mean the direct-product additive group with the $D$ coordinate more
significant in the order. Thus
\begin{equation}
 (d,n)<(e,m)
 \quad\Longleftrightarrow\quad
 d<e\ \text{ or }\ (d=e\text{ and }n<m).
 \label{eq:lex}
\end{equation}
Write $\one=(0,1)$ and let
\begin{equation}
 \MD=(\GD)_{\geq0}
      =(D_{>0}\times\Z)\cup(\{0\}\times\N).
 \label{eq:cone}
\end{equation}
The ambient product topology, whenever specified, is separate from
\eqref{eq:cone}'s arithmetic order.

\begin{theorem}[Polish and effective Presburger models]
\label{thm:main}
The following assertions hold.
\begin{enumerate}[label=\textup{(\arabic*)}]
\item With $D=\R$ in its usual topology, $\MR$ is an uncountable,
locally compact Polish space without isolated points. Coordinatewise
addition is continuous, and
\[
 (\MR;0,\one,+,<)\models\Th(\N;0,1,+,<).
\]
\item With $D=\Q^{\N}$, ordered by the first differing coordinate and
with each rational coordinate discrete, the space $\MB=\MD$ is
homeomorphic to $\Baire$. There is an explicit such homeomorphism under
which addition, successor, and division by each positive standard integer
are Type-2 computable.
\item In both models, every first-order definable relation, with
parameters, is $\Dtwo$. The order is neither open nor closed. In the Baire
presentation, each formula has a computable bound on the number of
changes in a convergent truth approximation, relative to its parameter
codes.
\item Both models embed as ordered additive submonoids of $\Oz$, fixing
$0$ and $1$. Neither admits any unital semiring multiplication with its
given addition, zero, and unit.
\end{enumerate}
\end{theorem}

The proofs occupy \cref{sec:algebra,sec:topology,sec:baire,sec:borel,sec:omnific,sec:multiplication}.
The first assertion alone answers the source question. The later assertions
explain how much regularity is possible and where it necessarily stops.

For orientation, an element $(r,-17)$ with $r>0$ is a \emph{positive}
element of $\MR$. Indeed it exceeds every standard numeral $(0,k)$.
Nevertheless,
\begin{equation}
 (1/j,0)\longrightarrow(0,0)
 \qquad(j\longrightarrow\infty).
 \label{eq:smallseq}
\end{equation}
in the chosen topology. Arithmetic size and topological proximity are
different notions. The construction would fail if one silently demanded
that this topology be the order topology.

\section{Full Presburger arithmetic, not merely a few axioms}
\label{sec:algebra}

\subsection{Discretely ordered groups with standard residues}

\begin{definition}
A \emph{$\Z$-group} in this article is a linearly ordered abelian group
$G$ with a distinguished least positive element $\one$, such that for
every positive standard integer $m$ and every $x\in G$ there exist
$q\in G$ and $r\in\{0,\ldots,m-1\}$ satisfying
\begin{equation}
 x=mq+r\one.
 \label{eq:division}
\end{equation}
Here $mq$ denotes repeated addition by a \emph{standard} integer, not a
new multiplication operation on $G$.
\end{definition}

Ordered abelian groups are torsion-free. The quotient and remainder in
\eqref{eq:division} are unique: if $m(q-q')=(r'-r)\one$ and $q>q'$, then
$q-q'\geq\one$, whereas $r'-r<m$; the other inequality is similar.
We suppress $\one$ in standard integer constants when no confusion is
possible.

\begin{proposition}[The split lexicographic group]
\label[proposition]{prop:zgroup}
If $D$ is a divisible ordered abelian group, then $D\lex\Z$ is a
$\Z$-group. Its standard division algorithm is
\begin{equation}
 (d,n)=m\bigl(d/m,\lfloor n/m\rfloor\bigr)+(0,n\bmod m).
 \label{eq:explicitdivision}
\end{equation}
If $(d,n)\geq0$, the displayed quotient is nonnegative.
\end{proposition}
\begin{proof}
Coordinatewise addition preserves the lexicographic order. A positive
$(d,n)$ either has $d>0$, in which case it exceeds $(0,1)$, or has $d=0$
and $n\geq1$. Thus $(0,1)$ is the least positive element. Divisibility
and torsion-freeness of $D$ give the unique element $d/m$. Ordinary
Euclidean division of $n$ now proves \eqref{eq:explicitdivision}.
If $d>0$, then $d/m>0$, regardless of the integer quotient. If $d=0$
and $n\geq0$, the integer quotient is nonnegative.
\end{proof}

The implication from these axioms to \emph{full} Presburger arithmetic is
important enough to prove. We use the classical elimination architecture
of Presburger and Cooper; \cite{Haase} supplies historical and algorithmic
context. The proof below makes explicit which algebraic hypotheses are
needed in a nonstandard group.

\subsection{Coset invariance versus standard-shift periodicity}
\label{sec:periodicity}

For $m\geq1$, write $mG=\{mh:h\in G\}$. A predicate $P$ is
\emph{$mG$-invariant} if
\begin{equation}
 y-z\in mG\quad\Longrightarrow\quad(P(y)\leftrightarrow P(z)).
 \label{eq:cosetinv}
\end{equation}
This is stronger than invariance under the single standard shift $m\one$.

\begin{example}[A false generalization of an integer lemma]
\label[example]{ex:falseperiod}
In $G=\R\lex\Z$, let $P(d,n)$ mean $d>0$. For every positive standard
$m$,
\[
 P(x+m\one)\leftrightarrow P(x)
\]
holds. Yet $P$ has witnesses and is false at every standard residue
$0,1,\ldots,m-1$. Thus
\[
 (\exists x\ P(x))\ \not\Longleftrightarrow\
 \bigvee_{0\leq r<m}P(r\one).
\]
It also fails the naive finite-interval reduction with lower bound $0$
and upper bound $(1,0)$. No flaw in the repository's \texttt{Int} theorem
is involved: its domain is the ordinary integers, where every multiple
is obtained by a standard finite iteration.
\end{example}

\begin{lemma}[The coset-invariant finite search]
\label[lemma]{lem:finite}
Let $G$ be a $\Z$-group, $m\geq1$ standard, and $P$ an $mG$-invariant
predicate. For $L,U\in G$,
\begin{align}
 &\exists y\ (L\leq y\leq U\ \wedge\ P(y))
 \nonumber\\[-.2em]
 &\hspace{1.3cm}\Longleftrightarrow
 \bigvee_{0\leq r<m}\bigl(L+r\leq U\ \wedge\ P(L+r)\bigr).
 \label{eq:finitesearch}
\end{align}
With only an upper bound, candidates $U-r$ suffice; with only a lower
bound, candidates $L+r$ suffice. With neither bound, the candidates are
$0,\ldots,m-1$.
\end{lemma}
\begin{proof}
Given a witness $y\geq L$, write $y-L=mh+r$ with $0\leq r<m$.
One has $h\geq0$: otherwise $h\leq-1$, and then
$mh+r\leq-m+(m-1)=-1$, contradicting $y-L\geq0$.
Consequently
\[
 L\leq L+r\leq y\leq U.
\]
The difference $y-(L+r)=mh$ lies in $mG$, so $P(L+r)$ holds.
Conversely every displayed candidate is already a witness. The same
argument applied to $U-y$ proves the upper-bound version. Without bounds,
use the standard remainder of $y$ itself.
\end{proof}

Congruence predicates have exactly the stronger invariance
\eqref{eq:cosetinv}. If $M$ is a common multiple of finitely many standard
moduli, any Boolean combination of congruences with those moduli is
$MG$-invariant. This remains true with arbitrary nonstandard parameters:
adding $Mh$ to a variable changes an affine congruence term by an element
of each relevant modulus subgroup.

\subsection{A self-contained elimination argument}

Expand the ordered-group language by predicates
\[
 D_m(t)\quad\Longleftrightarrow\quad t\in mG
 \qquad(m\geq2),
\]
and regard $D_1$ as always true. All of these are definable by existential
formulas in the group language.

\begin{theorem}[Classical Presburger elimination over $\Z$-groups]
\label{thm:qe}
In the expanded language
\[
 \{0,\one,+,-,\leq,(D_m)_{m\geq2}\},
\]
the axioms of $\Z$-groups eliminate quantifiers and form a complete theory.
In particular every $\Z$-group is elementarily equivalent to
$(\Z;0,1,+,-,\leq,(m\mid\cdot)_m)$.
\end{theorem}
\begin{proof}
It suffices to eliminate one existential quantifier from a quantifier-free
formula. Rewrite the formula in disjunctive normal form. Negated
inequalities are replaced using discreteness, for example
$\neg(s\leq t)\leftrightarrow t+1\leq s$. An equality is the conjunction
of two inequalities; a disequality is a disjunction of strict inequalities.
Finally, standard residue exhaustiveness gives
\[
 \neg D_m(t)\ \longleftrightarrow\
 \bigvee_{r=1}^{m-1} D_m(t-r).
\]
Thus it is enough to eliminate $\exists x$ from a conjunction of affine
inequalities and positive divisibility atoms, along with conjuncts not
involving $x$.

Let $A$ be a positive common multiple of every nonzero absolute coefficient
of $x$ occurring in that conjunction; take $A=1$ if there are none. Set
$y=Ax$ and add the condition $D_A(y)$. An atom
\[
 ax+t\leq0,\qquad a\neq0,
\]
can be multiplied by the positive standard integer $k=A/|a|$, giving
\[
 \operatorname{sgn}(a)y+kt\leq0.
\]
It is now an upper or lower bound on $y$. For a congruence, torsion-free
cancellation gives the equivalence
\begin{equation}
 D_m(ax+t)
 \quad\Longleftrightarrow\quad
 D_{km}\bigl(\operatorname{sgn}(a)y+kt\bigr).
 \label{eq:congscale}
\end{equation}
Indeed $kz\in kmG$ if and only if $z\in mG$: one direction multiplies a
witness, and the other cancels $k$ from $kz=kmh$.
The additional atom $D_A(y)$ is exactly what ensures that every candidate
$y$ comes from a unique $x$.

After normalizing signs, we have finitely many lower bounds $L_i\leq y$,
finitely many upper bounds $y\leq U_j$, and a conjunction $P(y)$ of
congruence conditions. Let $M$ be a positive common multiple of all their
moduli, including $A$. Then $P$ is $MG$-invariant.

If lower bounds exist, choose their maximum $L$. By
\cref{lem:finite}, a witness exists exactly when one of
$L,L+1,\ldots,L+M-1$ satisfies all the normalized conjuncts. The maximum
need not be introduced as a new function symbol: take a finite disjunction
over the indices $i$, require $L_h\leq L_i$ for every $h$, and substitute
$L_i+r$ for $y$. This produces a quantifier-free formula in the parameters.
If there are upper bounds but no lower bounds, choose a minimal $U_j$
and substitute $U_j-r$. If no bounds exist, substitute the standard
residues $r$. Conjuncts independent of $x$ are retained throughout.

This eliminates one existential quantifier. Induction on formula
complexity, replacing universal quantification by negated existential
quantification, eliminates all quantifiers.

A quantifier-free sentence uses only standard integer multiples of
$\one$. Their order and their residues are the same in every $\Z$-group,
by least positivity and uniqueness of standard remainders. Therefore all
$\Z$-groups satisfy the same sentences. The ordinary integers are one
such group, proving completeness and the final assertion.
\end{proof}

\begin{corollary}[The entire induction scheme]
\label[corollary]{cor:fullpresburger}
For every divisible ordered abelian group $D$,
\[
 (\MD;0,\one,+,<)\models\Th(\N;0,1,+,<).
\]
In particular it satisfies induction for every first-order formula of
this additive language, with arbitrary parameters from $\MD$.
\end{corollary}
\begin{proof}
By \cref{prop:zgroup,thm:qe}, $\GD\equiv\Z$. The nonnegative cone is
uniformly definable by $x\geq0$. Relativizing every quantifier to that
cone translates any sentence about $\MD$ into a sentence about $\GD$.
The same translation in $\Z$ gives the ordinary nonnegative integers.
Each universally closed induction instance is therefore true in $\MD$.
\end{proof}

\begin{proposition}[Definable least elements]
\label[proposition]{prop:least}
Every nonempty definable subset of $\MD$, allowing parameters, has a
least element.
\end{proposition}
\begin{proof}
Eliminate quantifiers and take a finite disjunction of conjunctions as
in \cref{thm:qe}. Include $x\geq0$ explicitly. For each conjunction,
the substitution $y=Ax$ introduces a lower bound $y\geq0$. Hence a
nonempty conjunction has a least solution among the finitely many valid
candidates at its maximum lower bound; the residue argument shows that
one of these candidates lies below every witness. Dividing by $A$ is
order-preserving. Finally take the minimum of the least solutions of the
nonempty disjuncts.
\end{proof}

This also gives a direct induction check after elimination: a nonempty
set of counterexamples to an induction instance would have a least
member. It cannot be $0$, and its predecessor would contradict the
induction step. The external standard part
$\{(0,n):n\in\N\}$ is a proper inductive subset when $D\neq0$, but is
not definable, even with parameters. Induction quantifies over definable
sets, not over every external subset.

\section{Polishness and a locally compact example}
\label{sec:topology}

A Polish space is a separable space whose topology is induced by some
complete metric. A particular convenient metric may be incomplete
without the space ceasing to be Polish. This distinction matters for
open rays.

\subsection{A useful exact criterion}

We first record the standard subspace criterion with a proof, so that the
completeness issue is visible.

\begin{lemma}[Polish subspaces]
\label[lemma]{lem:gdelta}
A subspace $A$ of a Polish space $X$ is Polish if and only if $A$ is a
$G_\delta$ subset of $X$.
\end{lemma}
\begin{proof}
Fix a compatible complete metric $d$ on $X$. If $A=\bigcap_j U_j$ with
$U_j$ open, consider on $A$ the continuous functions
$f_j(x)=1/d(x,X\setminus U_j)$; when $U_j=X$, set $f_j=0$.
The graph of $x\mapsto(f_j(x))_j$ is closed in $X\times\R^{\N}$.
Indeed a limit outside some $U_j$ would force $f_j$ to tend to infinity,
not to a real number. Projection identifies this closed graph
homeomorphically with $A$. Countable products and closed subspaces of
complete separable metric spaces are complete and separable.

Conversely, let $\delta$ be a complete metric inducing the subspace
topology on $A$. For each $n$, let $W_n$ be the union of all ambient open
sets $U$ that meet $A$ and satisfy
\[
 \operatorname{diam}_d(U)<2^{-n},\qquad
 \operatorname{diam}_{\delta}(U\cap A)<2^{-n}.
\]
These unions cover $A$. If $x\in\bigcap_n W_n$, choose $x\in U_n$ and
$a_n\in U_n\cap A$. Then $a_n\to x$ in $d$. For each fixed $n$, all
sufficiently late $a_k$ belong to $U_n$, so $(a_k)$ is $\delta$-Cauchy.
Its $\delta$-limit lies in $A$, and agreement of the two topologies
identifies that limit with $x$. Thus $A=\bigcap_n W_n$.
\end{proof}

\begin{theorem}[The positive-cone criterion]
\label{thm:criterion}
Let $D$ be a divisible ordered abelian group equipped with a Polish
topological-group topology, not necessarily its order topology. Give
$\Z$ the discrete topology. Then the subspace $\MD\subseteq D\times\Z$
is Polish if and only if $D_{>0}$ is $G_\delta$ in $D$. When this holds,
$D_{>0}$ is in fact $\Dtwo$, and addition on $\MD$ is continuous.
\end{theorem}
\begin{proof}
If $D_{>0}$ is $G_\delta$, so is its union with the closed singleton
$\{0\}$. On each clopen integer slice, the carrier \eqref{eq:cone} is
either $D_{>0}$ or $D_{\geq0}$. A countable topological sum of
$G_\delta$ subspaces is $G_\delta$ in the corresponding sum: its
complement is a countable union of $F_\sigma$ subsets of the clopen
slices. Apply \cref{lem:gdelta}.

Conversely, the slice with integer coordinate $-1$ is closed in $\MD$
and homeomorphic to $D_{>0}$. If $\MD$ is Polish, so is that slice, and
\cref{lem:gdelta} makes $D_{>0}$ a $G_\delta$ subset of $D$.
Negation is a homeomorphism on $D$, so $D_{<0}$ is $G_\delta$ as well.
The complement of $D_{>0}$ is $D_{<0}\cup\{0\}$, also $G_\delta$.
Thus $D_{>0}$ is both $F_\sigma$ and $G_\delta$.

Finally the cone is closed under coordinatewise addition. Restricting
continuous addition on $D\times\Z$ to a subspace with its image in that
subspace gives a continuous map $\MD^2\to\MD$.
\end{proof}

\subsection{The real-coordinate model, including a complete metric}

For $D=\R$, the positive cone is open. More explicitly, the slice
\[
 X_n=\{r:(r,n)\in\MR\}
\]
is $[0,\infty)$ when $n\geq0$ and $(0,\infty)$ when $n<0$. Hence
\begin{equation}
 \MR\ \cong\
 \coprod_{n\geq0}[0,\infty)\ \sqcup\
 \coprod_{n<0}(0,\infty).
 \label{eq:rays}
\end{equation}
This proves separability, local compactness, and absence of isolated
points directly.

For a fully explicit compatible complete metric, define
\[
 d_n(r,s)=
 \begin{cases}
 \min\{1,|r-s|\},&n\geq0,\\
 \min\{1,|\log r-\log s|\},&n<0,
 \end{cases}
\]
and set
\begin{equation}
 d((r,n),(s,m))=
 \begin{cases}d_n(r,s),&n=m,\\2,&n\neq m.\end{cases}
 \label{eq:completemetric}
\end{equation}
Every Cauchy sequence is eventually contained in one slice. The closed
ray with the first metric is complete; the open ray with the logarithmic
metric is isometric, before truncation, to $\R$. Thus \eqref{eq:completemetric}
is complete. It induces exactly the topology in \eqref{eq:rays}, which
is the original product-subspace topology. Continuity of addition follows
from that topology, not from a claim of uniform continuity in this metric.

The map $r\mapsto(r,0)$ on $r>0$ establishes uncountability.
Together with \cref{cor:fullpresburger}, this completes the simplest proof
of the affirmative answer.

\begin{remark}[Why the arithmetic order topology is not used]
The least positive element is $\one$. In $\GD$, the order interval
$(x-\one,x+\one)$ is the singleton $\{x\}$. The positive cone has the
same discreteness at its minimum. Consequently its arithmetic order
topology is discrete, and an uncountable such topology cannot be Polish.
The topology of \eqref{eq:rays} is different, as \eqref{eq:completemetric}
and the convergence in \eqref{eq:smallseq} explicitly demonstrate.
\end{remark}

\section{An explicit Baire-space realization}
\label{sec:baire}

Let
\[
 D_{\mathrm B}=\Q^{\N}
\]
with coordinatewise addition. Each factor $\Q$ is discrete, \emph{not}
given its Euclidean subspace topology. Two distinct sequences are compared
at the least index at which they differ. This lexicographic order makes
$D_{\mathrm B}$ a divisible ordered abelian group. Its product topology is
Polish and zero-dimensional; for example,
\[
 d_D(a,b)=2^{-\min\{j:a_j\neq b_j\}}
 \quad(a\neq b),\qquad d_D(a,a)=0
\]
is a compatible complete metric.

The sets
\[
 C_{k,q}=\{a:a_0=\cdots=a_{k-1}=0,\ a_k=q\},
 \qquad k\in\N,\ q\in\Q_{>0},
\]
are clopen cylinders, and
\begin{equation}
 (D_{\mathrm B})_{>0}=\coprod_{k,q>0}C_{k,q}.
 \label{eq:positivecylinders}
\end{equation}
No restriction is placed on later coordinates. Each cylinder is
homeomorphic to $\Q^{\N}$. The nonnegative cone
$(D_{\mathrm B})_{\geq0}$ is closed, since the negative cone is open.
\cref{thm:criterion} already proves that
\[
 \MB=\bigl(D_{\mathrm B}\lex\Z\bigr)_{\geq0}
\]
is Polish. We now identify its homeomorphism type without relying on a
space-classification theorem.

\subsection{The zero branch and its explicit tree code}

Fix computable bijections
\[
 e:\N\to\Q,\qquad e_+:\N\to\Q_{>0},\qquad
 e_{\geq0}:\N\to\Q_{\geq0},
\]
with computable inverses. A code $s\in\Baire$ determines a sequence
$a\in(D_{\mathrm B})_{\geq0}$ as follows. As long as all previously
chosen coefficients are zero, set $a_j=e_{\geq0}(s_j)$. Once the first
positive coefficient has appeared, use $a_j=e(s_j)$ at all later stages.
This is a bijection: the inverse uses the corresponding enumeration at
each step. A finite output prefix depends on exactly a finite input
prefix, in both directions. Thus it is a homeomorphism
\begin{equation}
 \Baire\ \cong\ (D_{\mathrm B})_{\geq0}.
 \label{eq:zerobranch}
\end{equation}
The all-zero rational sequence is not omitted; it is the branch that
continues choosing zero indefinitely.

\begin{theorem}[Baire-space carrier]
\label{thm:baire}
There is an explicit homeomorphism $h:\Baire\to\MB$.
In particular $\MB$ is uncountable, has no isolated points, and is
nowhere locally compact.
\end{theorem}
\begin{proof}
Partition $\MB$ into the following countably many clopen blocks. For each
$n\geq0$, take the entire slice
$(D_{\mathrm B})_{\geq0}\times\{n\}$. For $n<0$, partition that slice
further into $C_{k,q}\times\{n\}$ using
\eqref{eq:positivecylinders}. By \eqref{eq:zerobranch}, every nonnegative
slice is homeomorphic to Baire space; every negative-slice block is so
by coding its unrestricted rational tail. A countable disjoint sum of
copies of $\Baire$ is homeomorphic to $\Baire$: use the first natural
coordinate to choose a block and the remaining coordinates to choose a
point in it.

For completeness, no nonempty basic cylinder of $\Baire$ is relatively
compact. Its projection to the next unrestricted discrete coordinate
has infinite image, whereas the continuous image of a compact set in a
discrete space is finite. Every neighborhood contains such a cylinder.
The other two assertions follow from the elementary properties of the
product $\N^{\N}$.
\end{proof}

\subsection{A total code for every sequence of natural numbers}

Here is a fully specified block indexing, also used in the accompanying
program. Let $\langle a,b\rangle$ be Cantor's pairing function.
The first digit of a Baire code is a tag. An even tag $2n$ specifies the
nonnegative integer coordinate $n$ and the tree code
\eqref{eq:zerobranch}. An odd tag
\begin{equation}
 2\langle a,\langle k,b\rangle\rangle+1
 \label{eq:tag}
\end{equation}
specifies integer coordinate $-a-1$, first positive rational coordinate
$k$, and value $e_+(b)$ there. The remaining digits code the unrestricted
rational tail using $e$.

Every natural tag is of exactly one of these forms. Every infinite
natural sequence therefore codes a valid element; no hidden domain test
is required. Conversely, a point with negative integer coordinate has a
first positive rational coefficient, so its tag and tail are uniquely
recoverable. This describes $h$ and $h^{-1}$ explicitly.

\begin{theorem}[Computable continuous arithmetic]
\label{thm:effectiveops}
Under this presentation, the operations
\[
 x+y,\qquad x+\one,\qquad
 q_m(x)=\lfloor x/m\rfloor,\qquad r_m(x)=x\bmod m
\]
are Type-2 computable for every positive standard integer $m$. Here
$q_m(x)$ is the unique quotient in the model's Euclidean division by $m$,
not a real-valued floor operation. The remainders have their finite
discrete topology.
\end{theorem}
\begin{proof}
A Type-2 computation reads finitely many input digits to produce each
requested output digit. Decode the integer coordinates $n,z$ of the two
inputs. Each rational coefficient $a_j+b_j$ is computable from finite
prefixes. If $n+z\geq0$, output tag $2(n+z)$ and encode the coefficient
sequence by the nonnegative-cone tree. The tree decisions only inspect
finite prefixes.

If $n+z<0$, at least one input has a negative integer coordinate and hence
a strictly positive coefficient sequence. Its tag supplies an index $k$
of a positive leading coefficient. Since both coefficient sequences are
lex-nonnegative, their sum is strictly positive and its first nonzero
coordinate occurs at or before $k$. Search these finitely many
coordinates, output the negative tag \eqref{eq:tag}, and then encode the
tail. Thus the tag itself is computable; the algorithm does not attempt
to decide whether an arbitrary sequence is identically zero.

For division use
\[
 q_m(a,n)=(a/m,\lfloor n/m\rfloor),\qquad
 r_m(a,n)=n\bmod m.
\]
If the quotient's integer coordinate is negative, the input integer
coordinate was negative as well, and the same known leading index can
be used. Otherwise the nonnegative tree code applies. Constants and
successor are special cases of these finite-prefix procedures.
\end{proof}

The adjective ``computable'' here concerns operations on infinite
sequence codes. It does not assert that every element is computable,
or that equality and order of arbitrary streams can be decided in finite
time. The distinction will be quantified next.

\section{Definability, effective truth, and the optimal order rank}
\label{sec:borel}

A subset of a metrizable space is $\Dtwo$ when it is both $F_\sigma$ and
$G_\delta$. Finite Boolean combinations of such sets are again $\Dtwo$.
Closed sets, open sets, and clopen sets belong to this class.

\begin{theorem}[All definable relations have low Borel rank]
\label{thm:definableborel}
Under the hypotheses of \cref{thm:criterion}, every first-order definable
subset of $\MD^k$, with parameters, is $\Dtwo$.
\end{theorem}
\begin{proof}
In $\GD=D\times\Z$, the strict positive cone is
\[
 (D_{>0}\times\Z)\cup(\{0\}\times\Z_{>0}).
\]
It is $\Dtwo$, as is its complement. Affine group terms are continuous
maps to $\GD$, so their inequalities have $\Dtwo$ truth sets. Equalities
have closed truth sets because the ambient group is Hausdorff.

Divisibility is especially simple. Since $D$ is divisible,
\begin{equation}
 m\GD=D\times m\Z.
 \label{eq:modslice}
\end{equation}
This is clopen. Thus congruence atoms have clopen truth sets, including
when their affine terms contain parameters. Translate a cone formula
into the group and apply \cref{thm:qe}. The result is a finite Boolean
combination of these atoms. Finally restrict to $\MD^k$.
\end{proof}

\subsection{A formula-by-formula bound on changes of mind}

\begin{theorem}[Effective finite-mind-change truth]
\label{thm:mindchange}
For each first-order Presburger formula $\varphi(\bar x;\bar p)$, one can
compute a natural number $K_{\varphi}$ and an algorithm producing
\[
 A_s(\bar x;\bar p)\in\{0,1\}\qquad(s\in\N)
\]
from finite prefixes of the Baire codes, such that
\[
 \lim_s A_s(\bar x;\bar p)
  =\mathbf 1_{\MB\models\varphi(\bar x;\bar p)}.
\]
For every input tuple the displayed sequence changes value at most
$K_{\varphi}$ times. Parameter codes are supplied as oracles; they are
not assumed computable.
\end{theorem}
\begin{proof}
Compute a quantifier-free group formula equivalent to the relativized
cone formula, by \cref{thm:qe}. An affine group term has the form
$(c,n)$, where $n$ is an exactly computable integer and each rational
coefficient $c_j$ is computable from finite prefixes of the inputs.

To approximate the sign of $(c,n)$, first use the sign of $n$. At stage
$s$, inspect $c_0,\ldots,c_{s-1}$. If a nonzero coefficient has been
found, use the sign of the first such coefficient instead. If $c=0$,
the initial answer was correct. Otherwise the answer becomes correct
when the first nonzero coefficient is reached and never changes again.
An equality atom can be handled in the same way: it starts with the
answer determined by $n$ and becomes permanently false if a nonzero
coefficient is found. A congruence atom is decided immediately from $n$
by \eqref{eq:modslice}.

Evaluate the fixed Boolean combination of these approximations at each
stage. Its value can change only when an equality or inequality atom
changes. The number of those atom occurrences is a computable bound
$K_{\varphi}$. Each stage reads only finitely many digits. The truth set
is effectively $F_\sigma$ because it is
\[
 \bigcup_N\bigcap_{s\geq N}\{\bar x:A_s(\bar x;\bar p)=1\},
\]
where the stage sets are effectively clopen relative to the parameter
codes. The same expression with $0$ proves effective $G_\delta$ status.
\end{proof}

No effective stopping time is asserted: a later rational coordinate may
be the first nonzero one. Deciding sentences and deciding arbitrary
parameter instances are different problems. The former is ordinary
Presburger decidability; the latter is represented here by a convergent,
uniformly bounded-change process.

\subsection{Why rank one is impossible for the order}

\begin{theorem}[Open or closed arithmetic order forces discreteness]
\label{thm:orderoptimal}
Let $X$ carry a model of Presburger arithmetic and a topology in which
successor $x\mapsto x+\one$ is continuous. If either of the relations
$<$ or $\leq$ is open or closed in $X^2$, then $X$ is discrete.
Consequently, in any uncountable Polish model with continuous addition,
the arithmetic order is neither open nor closed.
\end{theorem}
\begin{proof}
If $\leq$ is open, then for fixed $a$ both sets
$\{x:x\leq a\}$ and $\{x:a\leq x\}$ are open. Their intersection is
$\{a\}$, so every point is isolated.

If $<$ is closed, taking the complement and reversing coordinates makes
$\leq$ open. If $<$ is open, the discrete-order equivalence
\[
 x\leq y\quad\Longleftrightarrow\quad x<y+\one
\]
and continuity of successor make $\leq$ open. Finally, if $\leq$ is
closed, then
\[
 x<y\quad\Longleftrightarrow\quad x+\one\leq y
\]
makes $<$ closed, reducing to the previous case.
An uncountable discrete space is not separable, hence not Polish.
\end{proof}

Together, \cref{thm:definableborel,thm:orderoptimal} make the $\Dtwo$
order bound optimal. This is a statement about the order relation, not
about every individual definable set: many congruence sets are clopen.
It also rules out a continuous Boolean-valued order test on the Baire
model, despite its computable addition.

\section{Which definable operations are continuous?}
\label{sec:continuity}

The coordinate topology makes several arithmetic operations exceptionally
regular, but not every definable operation. This section identifies exact
boundaries in the two explicit models.

\begin{proposition}[Euclidean self-similarity]
\label[proposition]{prop:selfsimilar}
In $\MR$ and $\MB$, for each positive standard integer $m$ the map
\begin{equation}
 x\longmapsto\bigl(q_m(x),r_m(x)\bigr)
 \label{eq:selfsimilar}
\end{equation}
is a homeomorphism from the model onto
$M\times\{0,\ldots,m-1\}$, where the finite factor is discrete.
Its inverse is $(q,r)\mapsto mq+r\one$.
\end{proposition}
\begin{proof}
Uniqueness and existence are \cref{prop:zgroup}. In both chosen
coefficient topologies, division of $d$ by $m$ is continuous. The integer
maps $n\mapsto\lfloor n/m\rfloor$ and $n\mapsto n\bmod m$ are
continuous on a discrete space. Hence \eqref{eq:selfsimilar} is
continuous into the product. Repeated addition and addition of a fixed
standard numeral make its inverse continuous.
\end{proof}

The same formula also explains why all standard congruences ignore the
entire coefficient sequence. Two elements have the same remainder modulo
every positive standard integer if and only if their integer coordinates
agree. The forward direction uses the elementary fact that an ordinary
integer divisible by every positive integer is zero. Their infinite
coordinates can still be arbitrary and different.

\begin{proposition}[The total predecessor has exactly one discontinuity]
\label[proposition]{prop:pred}
Define
\[
 p(0)=0,\qquad p(x)=x-\one\quad(x>0).
\]
In either explicit model, $p$ is continuous exactly at the points other
than $0$.
\end{proposition}
\begin{proof}
Away from $0$, this is a restriction of translation by $-\one$ in the
ambient topological group. Since $\{0\}$ is closed, the nonzero domain
is open, and translation has its image in the cone.

In the real model let $t_j=1/j$; in the Baire model let $t_j=e_j$, the
rational sequence with $1$ in coordinate $j$ and zeros elsewhere. In
either case $t_j>0$ and $t_j\to0$ in $D$. Then
\[
 (t_j,0)\to0,\qquad p(t_j,0)=(t_j,-1).
\]
The latter sequence does not converge to $p(0)=0$, because its integer
coordinate is always $-1$ and integer coordinates are discrete.
\end{proof}

\begin{theorem}[Exact discontinuity locus of truncated subtraction]
\label{thm:truncated}
For $M=\MR$ or $\MB$, let
\[
 T(x,y)=\max\{x-y,0\}
\]
using the arithmetic order in the ambient group. Its discontinuity locus
is exactly
\begin{equation}
 \bigl\{((d,n),(d,z))\in M^2:n\neq z\bigr\}.
 \label{eq:disclocus}
\end{equation}
\end{theorem}
\begin{proof}
Both positive and negative cones in the coefficient group $D$ are open
for our two examples. If the coefficient difference $d-e$ is nonzero,
its sign is locally fixed. Thus $T$ locally agrees either with the
continuous ambient difference or with the constant zero map.

Suppose instead that $x=y$. Along inputs sufficiently near this pair,
the integer difference is zero. The output is either their ambient
difference or zero. For any neighborhood $V$ of zero in the ambient
group, both possible outputs lie in $V$ when the inputs are sufficiently
near $(x,x)$. This proves continuity there.

It remains to consider equal coefficient coordinates and $k=n-z\neq0$.
Use the positive null sequence $t_j$ from \cref{prop:pred}.
If $k>0$, replace $y$ by $y+(t_j,0)$. These inputs converge to $(x,y)$,
but their truncated differences are zero, whereas $T(x,y)=(0,k)$.
If $k<0$, replace $x$ by $x+(t_j,0)$. The outputs are $(t_j,k)$, which
cannot converge to the original value $0$ because their integer
coordinate is the nonzero integer $k$. All perturbed inputs remain in
the positive cone. This proves exactly \eqref{eq:disclocus}.
\end{proof}

The standard numerals form a closed, relatively discrete subset of both
models, yet zero is not isolated in the whole space. This causes no
conflict with \cref{prop:pred}. Nor does continuity of standard division
imply continuity of every definable piecewise operation: the branch
boundary is decisive.

\section{Explicit set-sized fragments of the omnific integers}
\label{sec:omnific}

We use Conway normal-form notation
\[
 x=\sum_{\gamma\in S}r_\gamma\omega^\gamma,
\]
where $S$ is reverse well-ordered and all coefficients are real. In the
omnific ring the occurring exponents are nonnegative, and the coefficient
at exponent zero is an ordinary integer. This is the interface recorded
in the repository's surreal development \cite{RepoSurreal}. Addition is
coefficientwise, and a nonzero normal form has the sign of its leading
coefficient.

\begin{theorem}[One-scale and countable-scale realizations]
\label{thm:omnific}
There are injective ordered additive maps, fixing $0$ and $1$,
\begin{align}
 \iota_{\R}:\R\lex\Z&\longrightarrow\Oz,
 &\iota_{\R}(r,n)&=r\omega+n,
 \label{eq:realembedding}\\
 \iota_{\mathrm B}:\Q^{\N}\lex\Z&\longrightarrow\Oz,
 &\iota_{\mathrm B}(a,n)&=
       \sum_{j=0}^{\infty}a_j\omega^{1/(j+1)}+n.
 \label{eq:baireembedding}
\end{align}
Their positive images, equipped with the transported coefficient
topologies, are respectively the Polish models $\MR$ and $\MB$.
\end{theorem}
\begin{proof}
The first formula is a finite normal form. For the second, the fixed
exponents
\[
 1,\tfrac12,\tfrac13,\ldots
\]
strictly decrease and are positive. Every nonempty subset of these
exponents has a largest element, namely the one with least index.
Thus every possible support in \eqref{eq:baireembedding} is reverse
well-ordered, and adjoining exponent zero preserves that property.
The expression is a legitimate normal form, even when infinitely many
coefficients are nonzero. All positive-exponent coefficients are real,
and the constant term is the integer $n$, so the value is omnific.

Uniqueness of normal forms gives injectivity. Addition of two such
expressions is coordinatewise on a common fixed support, giving
additivity. The leading nonzero coefficient is the first nonzero $a_j$;
if all are zero, the sign is that of $n$. Therefore the map preserves
and reflects the lexicographic order. The same reasoning proves the
one-scale case. Transporting a topology through a bijection onto its
image preserves all the asserted topological properties.
\end{proof}

\subsection{What this realization does and does not preserve}

These are ordered additive embeddings, \emph{not} ring embeddings. Both
images contain $\omega$, but neither contains $\omega^2$. Their
multiplicative nonclosure is therefore immediate from the normal forms.
The stronger theorem in \cref{sec:multiplication} will show that replacing
surreal multiplication by some other multiplication cannot repair the
problem while retaining the same carrier, addition, and unit.

The topologies in \cref{thm:omnific} are precisely specified coefficient
topologies. The real coordinate is Euclidean in
\eqref{eq:realembedding}; the rational coordinates are discrete and
combined by a product topology in \eqref{eq:baireembedding}. Neither is
the discrete arithmetic order topology. No assertion that they are a
canonical topology on all omnific integers is intended.

The ordinary constant-term map becomes simply
\[
 \ct\bigl(\iota(d,n)\bigr)=n.
\]
It is continuous into the discrete integers for these models. Standard
congruence data see only this coordinate, matching
\eqref{eq:modslice}. The standard part itself remains externally
specified and nondefinable in the Presburger structure.

\subsection{Why no proper-class difficulty arises}

The source sets $\R\times\Z$ and $\Q^{\N}\times\Z$ are ordinary sets.
The normal-form constructions above are definable assignments on those
sets. Replacement therefore gives set-sized images. The set of birthdays
of their elements has an ordinal supremum; choosing a larger ordinal
places both images inside a set-sized birthday cutoff of $\No$.
No explicit optimal cutoff is needed for any theorem here.

Accordingly, one can conduct the entire argument in the product-group
models and use \cref{thm:omnific} only as a final realization theorem.
This is also the most economical formalization strategy. It avoids
making a class-sized object the carrier of a Polish space or a
first-order structure in ordinary set-based semantics.

\section{An absolute obstruction to multiplication}
\label{sec:multiplication}

A Borel obstruction would not rule out a wildly discontinuous
multiplication. For the models constructed here, the obstruction is
stronger and purely algebraic.

\begin{definition}
In a Presburger model $M$, a positive element $u$ is a
\emph{cofinal infinite additive unit} if
\[
 u>n\one\quad\text{for every standard }n,
\]
and for every $x\in M$ there is a positive standard integer $N$ with
$x\leq Nu$.
The word ``unit'' here refers only to cofinality under repeated addition;
it is not a multiplicative unit.
\end{definition}

\begin{theorem}[Cofinal-infinite-element obstruction]
\label{thm:no_semiring}
A Presburger model with a cofinal infinite additive unit cannot be
expanded to a unital semiring with the same addition, zero, and
multiplicative unit $\one$. No continuity, measurability, commutativity
of multiplication, or induction involving multiplication is assumed.
\end{theorem}
\begin{proof}
Suppose such a multiplication exists. Cofinality applied to $u^2$ gives
a positive standard $N$ such that $u^2\leq Nu$. Since $u$ exceeds every
standard numeral, it has a decomposition
\[
 u=(N+1)\one+v\qquad(v\in M).
\]
Distributivity and the unit law give
\[
 u^2=(N+1)u+vu\geq(N+1)u>Nu,
\]
a contradiction. The first inequality uses only that $vu$ belongs to the
nonnegative carrier. The last one follows from $u>0$ and additive
cancellation. In any semiring expansion of this additive monoid, its
natural order $x\leq y\leftrightarrow\exists z\ (x+z=y)$ is already the
Presburger order, so no additional order-compatibility axiom was hidden
in the argument.
\end{proof}

\begin{corollary}
\label[corollary]{cor:no_semiring_models}
Neither $\MR$ nor $\MB$ admits a unital semiring multiplication with the
displayed addition and unit.
\end{corollary}
\begin{proof}
In $\MR$ take $u=(1,0)$. It exceeds every standard numeral, and an
integer $N>r$ gives $(r,n)<Nu$. In $\MB$ take $u=(e_0,0)$, where $e_0$
has first coefficient $1$ and all later coefficients zero. For any
$(a,n)$ choose an integer $N>a_0$; comparison at coordinate zero gives
$(a,n)<Nu$. Apply \cref{thm:no_semiring}.
\end{proof}

In their omnific realizations, this cofinal element is $\omega$.
The missing larger scale $\omega^2$ is not merely a defect of one chosen
multiplication: every possible semiring multiplication would have to
produce an element growing faster than every standard multiple of
$\omega$, which the additive carrier cannot contain.

For a different uncountable Polish Presburger model that does allow a
semiring expansion algebraically, Glazer's cited Borel obstruction for
$Q+I\mathrm{Open}$ would still matter. Our argument is not a replacement
for that general theorem; it identifies a stronger obstruction in these
particular additive structures.

\section{Endomorphisms, elementary submodels, and a contrast}
\label{sec:embeddings}

The two models have markedly different self-embedding behavior. All
homomorphisms in the next theorem preserve addition, zero, and $\one$;
continuity and injectivity are \emph{not} assumed.

\begin{theorem}[Rigidity of the real-coordinate model]
\label{thm:realendos}
Every unital additive endomorphism $f:\MR\to\MR$ has the form
\begin{equation}
 f(r,n)=(cr,n)\qquad(c\in\R_{>0}).
 \label{eq:endo}
\end{equation}
Conversely each such formula is an endomorphism. All of them are
homeomorphisms and Presburger automorphisms. In particular $\MR$ has no
proper elementary self-embedding.
\end{theorem}
\begin{proof}
The group completion of the cancellative monoid $\MR$ is
$G=\R\lex\Z$. Thus $f$ extends uniquely to an additive group map
$F:G\to G$ fixing $\one$. Inside $G$,
\[
 \bigcap_{m\geq1}mG=\R\times\{0\}.
\]
One inclusion follows from real divisibility; the other follows because
the integer coordinate of an element of the intersection is divisible
by every positive integer. Every additive group map preserves this
intersection. Hence
\[
 F(r,0)=(h(r),0),\qquad F(0,n)=(0,n)
\]
for some additive map $h:\R\to\R$.

Since $f$ maps the positive cone into itself, $h(r)\geq0$ whenever
$r\geq0$. Thus $h$ is monotone. If $c=h(1)$, rational comparisons and
additivity give $h(r)=cr$ for every real $r$: sandwich $r$ between
rational numbers and use the corresponding inequalities for $h$.
Here $c\geq0$. But $c=0$ would send the positive element $(1,-1)$ to
$(0,-1)$, outside the cone. Thus $c>0$.

Formula \eqref{eq:endo} and its inverse with scalar $1/c$ preserve the
cone, addition, order, and unit, and are continuous in the product
coordinates. This proves all assertions.
\end{proof}

\begin{theorem}[All elementary submodels of the real model]
\label{thm:submodels}
The elementary submodels of $\MR$ are exactly
\begin{equation}
 M_E=(E\lex\Z)_{\geq0},
 \qquad E\text{ a }\Q\text{-linear subspace of }\R.
 \label{eq:submodels}
\end{equation}
Among them, the only ones Polish in the inherited topology are the
standard copy $\N$ and the whole model $\MR$. Every nonzero proper
choice of $E$ gives a dense, non-Polish elementary submodel.
\end{theorem}
\begin{proof}
If $E$ is a $\Q$-linear subspace, then $E\lex\Z$ is a $\Z$-group.
Its inclusion in $\R\lex\Z$ preserves and reflects the expanded
quantifier-free language, including all congruences, since a standard
modulus depends only on the integer coordinate. Quantifier elimination
makes this inclusion elementary. Relativizing to positive cones proves
$M_E\preccurlyeq\MR$.

Conversely, let $S\preccurlyeq\MR$, and let $H=S-S$ be its group
completion inside $G$. If $a,b\in S$ and $a-b\geq0$, the unique solution
of $b+x=a$ lies in $S$ by elementarity. Therefore
$S=H\cap\MR$. Since $H$ contains every standard integer, subtracting
the integer coordinate from any $(r,n)\in H$ shows that
$H=E\times\Z$ for an additive subgroup $E\leq\R$.
For $r\in E_{\geq0}$, the model's unique division of $(r,0)$ by a
positive standard $m$ has remainder zero and quotient $(r/m,0)$.
Elementarity puts that quotient in $S$. Applying the same observation
to negatives in $H$ shows $E$ is divisible, hence $\Q$-linear.

If $E\neq0$, it contains all rational multiples of a nonzero real and
is dense in $\R$. Thus $M_E$ is dense in every slice of $\MR$.
Suppose $M_E$ is also Polish in the inherited topology. Its closed slice
of integer coordinate zero is homeomorphic to $E\cap[0,\infty)$.
By \cref{lem:gdelta}, this is $G_\delta$ in $\R$. Its negative is also
$G_\delta$, so their union $E$ is $G_\delta$ in $\R$.
For any $r\in\R$, the dense $G_\delta$ sets $E$ and $r+E$ intersect
by the Baire category theorem. Hence $r\in E-E=E$, giving $E=\R$.
Finally $E=0$ gives the relatively discrete standard copy $\N$, which
is Polish, and $E=\R$ gives the whole model.
\end{proof}

For example, $E=\Q$ gives a countable dense elementary submodel. Its
inherited topology is not Polish; there is no inconsistency between its
density and the complete metric on the ambient model.

\begin{theorem}[Clopen elementary self-copies of the Baire model]
\label{thm:shift}
Let
\[
 \sigma((a_0,a_1,a_2,\ldots),n)
   =((0,a_0,a_1,a_2,\ldots),n).
\]
Then $\sigma:\MB\to\MB$ is a proper elementary embedding and a
homeomorphism onto a clopen subspace. Its iterates satisfy
\[
 \bigcap_{k\geq0}\sigma^k(\MB)=\{(0,n):n\in\N\}.
\]
Thus a descending sequence of proper clopen elementary self-copies has
intersection exactly the standard elementary submodel.
\end{theorem}
\begin{proof}
Inserting an initial zero preserves and reflects the lexicographic
order, addition, and every standard congruence. It maps the divisible
coefficient group into the subgroup with first coordinate zero.
Quantifier elimination gives elementarity, first for the groups and
then for their cones. The image is exactly the subset with $a_0=0$,
which is clopen because the first rational coordinate is discrete.
Deleting that zero is a continuous inverse on the image. It is proper,
since elements with $a_0>0$ are omitted.
The $k$-th image has its first $k$ rational coordinates zero. Membership
in all images therefore means that every rational coordinate is zero;
positivity then forces $n\geq0$.
\end{proof}

These clopen submodels are external subsets, not Presburger-definable
subsets, even with parameters: a proper definable subset containing zero
and closed under successor would violate induction. Topological
simplicity and internal definability are different properties.

\section{A compactness obstruction independent of multiplication}
\label{sec:compact}

The positive answer on Baire space does not extend to every familiar
uncountable Polish space. In particular compact carriers are impossible.
The underlying fact is the standard compact-cancellative-monoid argument,
which we include in its commutative form.

\begin{lemma}[Compact cancellation produces inverses]
\label[lemma]{lem:compactmonoid}
Every compact Hausdorff commutative cancellative topological monoid is
a group algebraically.
\end{lemma}
\begin{proof}
Write the identity as $0$ and fix $a$. Let $S$ be the closure of
$\{na:n\geq1\}$. Continuity of addition makes $S$ a nonempty compact
subsemigroup. By the compact intersection property and Zorn's lemma,
$S$ has a minimal nonempty compact subsemigroup $K$.

For $x\in K$, the set $x+K$ is a nonempty compact subset of $K$ and is a
subsemigroup: for $u,v\in K$,
\[
 (x+u)+(x+v)=x+(x+u+v)\in x+K.
\]
Minimality gives $x+K=K$. Since $x\in K$, there is $y\in K$ with
$x+y=x=x+0$, and cancellation gives $y=0$. Thus $0\in K\subseteq S$.
But $a+X$ is compact and hence closed, and contains every positive
multiple of $a$. Therefore it contains $S$ and in particular $0$.
So $a+b=0$ for some $b\in X$.
\end{proof}

\begin{corollary}[No compact positive-cone model]
\label[corollary]{cor:compact}
No model of Presburger arithmetic has a compact Hausdorff topology with
continuous addition. In particular Cantor space cannot be a carrier,
although Baire space can.
\end{corollary}
\begin{proof}
Its addition is a commutative cancellative monoid. The lemma would give
an additive inverse of $\one$, but no nonnegative element can sum with
$\one$ to zero.
\end{proof}

This obstruction uses neither Borel definability nor multiplication. It
is separate from \cref{thm:no_semiring}: the real model is locally
compact but not compact, and the Baire model is nowhere locally compact.

\section{Formalization and reproducibility}
\label{sec:formalization}

\subsection{A minimal theorem dependency graph}

The construction naturally splits into independently checkable layers.
The core algebra does not require real analysis or surreal numbers.
First define a discretely ordered abelian group with standard residue
representatives. Prove the coset-invariant finite-search lemma. Then
prove coefficient normalization, elimination, completeness, and the
positive-cone interpretation. Only afterward supply the real or
rational-sequence carrier and its topology. The omnific realization can
be verified last, as an independent bridge.

A practical Lean development can use the following \emph{proposed}
modules. These names are a plan, not existing checked files:
\begin{center}
\begin{tabular}{@{}>{\raggedright\arraybackslash}p{.31\textwidth}>{\raggedright\arraybackslash}p{.61\textwidth}@{}}
\toprule
Proposed module & Required mathematical content\\
\midrule
\texttt{ZGroup.Basic} & Least positive element, standard residue witnesses,
uniqueness, and order properties of the quotient.\\
\texttt{ZGroup.CosetCooper} & Full $mG$-invariance and the finite-search
criterion; the counterexample to weaker periodicity.\\
\texttt{ZGroup.Elimination} & Affine-term semantics over an arbitrary
$\Z$-group, coefficient scaling, quantifier elimination.\\
\texttt{PolishPresburger.}\allowbreak\texttt{Real} & The positive cone of $\R\lex\Z$ and
its complete slice metric.\\
\texttt{PolishPresburger.}\allowbreak\texttt{Baire} & The rational-sequence cone and the
explicit block/tree homeomorphism.\\
\texttt{PolishPresburger.}\allowbreak\texttt{Structure} & Order-rank obstruction, absence of
semiring expansions, and self-embedding theorems.\\
\texttt{PolishPresburger.}\allowbreak\texttt{Omnific} & The fixed-support normal-form maps and
their additive/order laws.\\
\bottomrule
\end{tabular}
\end{center}

\subsection{The nonstandard Cooper boundary must be explicit}

The inspected repository theorem \texttt{periodic\_interval} has an
ordinary-integer predicate and a one-step periodicity hypothesis
\cite{RepoCooper}. Its proof can iterate by any ordinary integer quotient.
Over a general $\Z$-group, the quotient in $x-L=mq+r$ is an element of
that group, not necessarily a standard numeral. A proof by ordinary
natural-number induction cannot supply invariance under that shift for
an arbitrary external predicate.

The correct interface is therefore a predicate on the quotient $G/mG$,
or equivalently the hypothesis \eqref{eq:cosetinv}. Affine congruence
atoms satisfy it by algebra, and Boolean combinations preserve it. This
makes the generalized arithmetic lemma reusable and prevents an
incorrect formal transfer from being hidden by notation. The original
integer theorem should be retained as an instance; it does not need to
be weakened or retracted.

The existing decision procedure's semantics are over integers. Reusing
its syntax and transformations is plausible, but proving soundness over
all $\Z$-groups is an additional obligation. The repository's existence
of an executable integer eliminator, by itself, does not certify
\cref{cor:fullpresburger} for a new carrier.

\subsection{Executable material included with this manuscript}

\texttt{polish\_presburger.py} implements exact finite-rank rational
arithmetic and the Baire codec of \cref{sec:baire}. The codec uses a
computable positive-rational enumeration and Cantor pairing, supplies
addition and standard division on infinite-sequence oracles, and gives
a one-change comparison approximation. Every Baire input code is valid.
The lower-level encoding routine takes a positive-cone promise, documented
in its interface; the public arithmetic operations preserve that promise.

\texttt{verify.py} performs deterministic tests using exact
\texttt{fractions.Fraction} arithmetic and a fixed seed. The executed
run produced \textbf{17,930 successful checks}. They cover pairing and
rational-enumeration round trips, additive and order identities,
Euclidean division, cone closure, coset-invariant finite candidates,
Baire-code prefix round trips, arithmetic output prefixes, the comparison
change bound, and the predecessor boundary witness. The deliberately
false weak-periodicity generalization is tested as a counterexample,
not asserted as a theorem.

\begin{tcolorbox}[title=Reproduction]
\raggedright
With Python 3.10 or later, run \texttt{python verify.py}. No external
Python package is needed. The script writes
\texttt{verification\_results.json}. Compile the article twice with
\texttt{pdflatex -interaction=nonstopmode -halt-on-error
polish\_presburger\_glazer.tex} to resolve references.
\end{tcolorbox}

Tests on finite-rank elements do not quantify over arbitrary elements of
$\Q^{\N}$, and finite prefixes do not decide equality of infinite
streams. Polishness, induction, and universal semantic statements are
established by the proofs, not by the test count. No new Lean or Rocq
formalization is included in this archive.

\section{Further research questions and topics}
\label{sec:questions}

The following questions are proposed directions arising from the proved
results. Except where explicitly linked to the source question, their
status is not presented as a certification that the literature has no
answer.

\begin{question}[A small formal certificate for the affirmative answer]
Can the one-scale model, the full $\Z$-group elimination theorem, and
its Polish topology be formalized with a compact audit boundary in
ProveIt? The key initial milestone is the generic coset-invariant Cooper
lemma, rather than merely a finite test of arithmetic identities.
\end{question}

\begin{question}[Prescribed zero-dimensional locally compact carriers]
Which uncountable zero-dimensional locally compact Polish spaces support
continuous-addition Presburger models? Compact spaces are excluded and
Baire space is realized, but Baire space is nowhere locally compact.
This leaves a meaningful topological region not settled by our examples.
\end{question}

\begin{question}[Connected and finite-dimensional carriers]
What additional obstruction prevents or permits a continuous-addition
model on a connected Polish space? Our locally compact model has
countably many ray components, and the arithmetic order cannot be open
or closed. A proof for a prescribed connected carrier must use more than
the order-topology observation alone.
\end{question}

\begin{question}[Removing the algebraic obstruction to multiplication]
Can one give an equally explicit Polish continuous-addition Presburger
model whose underlying additive structure admits a unital semiring
expansion? Any candidate must avoid a cofinal infinite additive unit.
If one also requires $Q+I\mathrm{Open}$, the multiplication cannot be
Borel by the cited obstruction. Separating algebraic expandability from
Borel expandability is the central issue.
\end{question}

\begin{question}[Classifying admissible coefficient groups]
Which divisible ordered Polish abelian groups have $G_\delta$ positive
cone? \Cref{thm:criterion} turns this into an exact criterion for the
split construction. Can their possible underlying topologies, order
ranks, and continuous order-preserving endomorphisms be classified under
useful additional hypotheses?
\end{question}

\begin{question}[Continuity of a richer definable language]
How much of Presburger definable choice can be made continuous in one
uncountable Polish presentation? Standard quotients and remainders are
continuous here, but total predecessor and truncated subtraction are
not everywhere continuous. Is there a natural maximal collection of
definable operations admitting simultaneous continuity?
\end{question}

\begin{question}[Optimal finite-mind-change complexity]
For a given Presburger formula, what is the smallest possible uniform
change bound in the Baire presentation? The proof gives the number of
order/equality atom occurrences after elimination. That bound need not
be optimal and may be much larger than necessary. Relating the minimal
bound to the finite difference hierarchy or to formula structure would
refine the $\Dtwo$ result.
\end{question}

\begin{question}[Polish elementary-submodel lattices]
Which lattices of closed or clopen elementary submodels can occur in a
Polish continuous-addition Presburger model? The real model has only two
Polish elementary submodels with inherited topology, whereas the Baire
model has infinitely many proper clopen elementary self-copies with
standard intersection. What intermediate behaviors can be realized?
\end{question}

\begin{question}[Canonical surreal cutoffs and topologies]
Find useful explicit birthday bounds for the two fixed-support omnific
images and formalize the coefficient-topology transport inside those
cutoffs. More generally, which support sets admit Polish coefficient
topologies with continuous addition and full Presburger induction on
the resulting positive cone? Such a classification should keep ring
closure separate from additive closure.
\end{question}

\begin{question}[Metatheoretic strength]
What weak foundational theory suffices to certify the affirmative
construction, including the chosen representation of Polishness?
The real model has an explicit complete metric, and the Baire model has
an explicit sequence code, so a general model-existence theorem is
unnecessary. This is distinct from Glazer's Question 1 concerning the
strength needed for his stronger-arithmetic obstruction.
\end{question}

\section{Conclusion}

The affirmative construction can be summarized by one change of viewpoint:
use the lexicographic order to interpret arithmetic, and a coefficient
topology to interpret continuity. Divisibility of the leading-coordinate
group supplies all standard residues; quantifier elimination then
supplies full Presburger induction. Neither countability of the standard
model nor discreteness of the arithmetic order forces the external
presentation to be countable or discrete.

The real model gives a locally compact answer, while the rational-sequence
model gives a genuinely effective answer on Baire space. Their order
relations have the smallest possible nontrivial Borel rank, yet they
exhibit different elementary self-embedding behavior. Their realization
inside $\Oz$ connects these results directly to the repository's
normal-form arithmetic. Finally, the cofinal-infinite-element argument
shows exactly why these additive examples cannot be mistaken for
topological models of full arithmetic.

\appendix
\section{Worked examples and an audit of the hypotheses}
\label{sec:audit}

\subsection{A negative constant term is compatible with positivity}

In $\MR$, take
\[
 x=(\tfrac32,-17),\qquad y=(\tfrac14,6).
\]
Both are positive, and
\[
 x+y=(\tfrac74,-11).
\]
Division of $x$ by $5$ gives
\[
 q_5(x)=(\tfrac3{10},-4),\qquad r_5(x)=3,
\]
because $-17=5(-4)+3$. The quotient is positive because its real
coordinate is positive. A definition that incorrectly demanded a
nonnegative integer coordinate for every element would break the
predecessor and division axioms.

\subsection{A nonstandard finite-candidate calculation}

Let $L=(\tfrac12,-5)$, $U=(\tfrac23,0)$, and require
\[
 y\equiv(\tfrac7{13},2)\pmod{3G}.
\]
The real coefficient of a congruence parameter is irrelevant because
$\R$ is divisible. Among $L,L+1,L+2$, the candidate $L+1$ has integer
coordinate $-4\equiv2\pmod3$, so it is a witness. Every candidate lies
below $U$, since its leading coefficient is $1/2<2/3$.
No iteration a nonstandard number of times is used: the finite search is
justified by full invariance on $3G$-cosets.

\subsection{A Baire convergence that order cannot detect continuously}

Let $e_j$ be the rational unit vector at coordinate $j$. Then
\[
 (e_j,0)\to(0,0),\qquad (e_j,0)>n\one
 \quad\text{for every fixed standard }n.
\]
The initial coefficients eventually agree with zero to every prescribed
finite length. In omnific notation this sequence is
\[
 \omega,\quad\omega^{1/2},\quad\omega^{1/3},\quad\ldots,
\]
which converges to zero in the \emph{specified coefficient topology},
not in the arithmetic order topology. This is a concrete illustration,
not a claim about an unspecified topology on surreal numbers.

\subsection{Hypothesis checklist}

\begin{longtable}{@{}>{\raggedright\arraybackslash}p{.30\textwidth}>{\raggedright\arraybackslash}p{.62\textwidth}@{}}
\toprule
Potential pitfall & Resolution in the proofs\\
\midrule
\endhead
Only the division axioms are checked &
\Cref{thm:qe,cor:fullpresburger,prop:least} establish the complete theory
and every additive induction instance.\\[.6em]
An open ray is given an incomplete metric &
\Cref{sec:topology} provides a compatible logarithmic metric on each
open ray and a complete metric on the whole sum.\\[.6em]
The intended topology is silently replaced by the order topology &
Both topologies are explicitly distinguished; the arithmetic order
topology is discrete.\\[.6em]
Standard-shift invariance is used with a nonstandard quotient &
\Cref{ex:falseperiod} exhibits the failure; \cref{lem:finite} uses
invariance on $mG$-cosets instead.\\[.6em]
The all-zero rational branch is lost &
The nonnegative-cone tree in \eqref{eq:zerobranch} explicitly includes it.\\[.6em]
Computing an output tag secretly requires deciding equality of streams &
For a negative output integer coordinate, a negative input supplies a
finite bound on the first positive coefficient.\\[.6em]
A type-theoretic ``computable model'' is confused with a countable one &
The presentation is Type-2 on all Baire sequences; not every element is
computable, and order has only a limit decision procedure.\\[.6em]
The omnific image is assumed closed under multiplication &
It is only an additive fragment. \Cref{thm:no_semiring} rules out even
an alternative semiring multiplication on the same carrier.\\[.6em]
A proper class is treated as a Polish space &
Only explicit set-sized images and a sufficiently large set-sized
birthday cutoff are used.\\[.6em]
Existing repository proofs are credited with new theorems &
The formalization roadmap distinguishes the inspected integer theorems
from the additional nonstandard and topological obligations.\\[.6em]
Numerical testing is mistaken for a proof &
The exact test scope is documented in \cref{sec:formalization}; no formal
verification claim is made for this article.\\[.6em]
An affirmative proof is mistaken for certified historical priority &
The source question and version are identified; independent priority and
refereeing remain separate from the mathematical proof.\\
\bottomrule
\end{longtable}

\clearpage
\begin{thebibliography}{9}
\bibitem{Glazer}
Elliot Glazer,
\emph{A Topological Tennenbaum Theorem},
arXiv:2311.13699v1, submitted November 22, 2023.
Question 2 is on page 8; the stronger-arithmetic obstruction used for
comparison is Corollary 1 on page 3.
\href{https://arxiv.org/abs/2311.13699}{arXiv record};
\href{https://arxiv.org/pdf/2311.13699v1}{versioned manuscript}.

\bibitem{EHW}
Ali Enayat, Joel David Hamkins, and Bartosz Wcis\l o,
\emph{Topological models of arithmetic},
arXiv:1808.01270; published in \emph{Fundamenta Mathematicae}
256 (2022), 171--193.
\href{https://arxiv.org/abs/1808.01270}{arXiv record}.

\bibitem{Haase}
Christoph Haase,
\emph{A survival guide to Presburger arithmetic},
\emph{ACM SIGLOG News} 5(3) (2018), 67--82.
DOI: \href{https://doi.org/10.1145/3242953.3242964}{10.1145/3242953.3242964}.
\href{https://www.cs.ox.ac.uk/people/christoph.haase/home/publication/haa-18/haa-18.pdf}{Author's manuscript}.

\bibitem{RepoPresburger}
Vladimir Reshetnikov, \emph{ProveIt: Presburger arithmetic is decidable},
\path{Logic/PresburgerArithmetic/README.md}, repository snapshot
\texttt{a11efab09a2d08f692a09829cb6d3440ddd1795a}.
\href{https://github.com/VladimirReshetnikov/ProveIt/blob/a11efab09a2d08f692a09829cb6d3440ddd1795a/Logic/PresburgerArithmetic/README.md}{Pinned repository file}.

\bibitem{RepoCooper}
Vladimir Reshetnikov, \emph{ProveIt: the arithmetic heart of Cooper
elimination},
\path{Logic/PresburgerArithmetic/Lean/PresburgerArithmetic/Cooper.lean},
lines 1--200 inspected at the same snapshot.
\href{https://github.com/VladimirReshetnikov/ProveIt/blob/a11efab09a2d08f692a09829cb6d3440ddd1795a/Logic/PresburgerArithmetic/Lean/PresburgerArithmetic/Cooper.lean}{Pinned source file}.

\bibitem{RepoSurreal}
Vladimir Reshetnikov, \emph{ProveIt: Surreal},
\path{Algebra/SurrealNumbers/README.md}, lines 1--230 inspected at the
same snapshot; see the formalization discussion of normal forms,
omnific integers, least positivity, and standard congruences.
\href{https://github.com/VladimirReshetnikov/ProveIt/blob/a11efab09a2d08f692a09829cb6d3440ddd1795a/Algebra/SurrealNumbers/README.md}{Pinned repository file}.
\end{thebibliography}

\end{document}
